% MyFeelsLike: a personalized apparent-temperature model
% Draft for arXiv (physics.ao-ph / stat.AP). Compiles with pdflatex, no
% non-standard packages.
\documentclass[11pt]{article}

\usepackage[margin=1in]{geometry}
\usepackage{amsmath,amssymb}
\usepackage{booktabs}
\usepackage{graphicx}
\usepackage{hyperref}
\usepackage{natbib}
% The reference list is inline (see thebibliography at the end) rather than
% generated by BibTeX. arXiv does not run BibTeX -- it expects a .bbl -- and a
% single local pdflatex pass would otherwise leave every citation as "?".
% myfeelslike.bib holds the same entries if you would rather generate it.

\title{MyFeelsLike: Learning an Individual Apparent-Temperature Scale\\
from Self-Reported Comfort Ratings}

% TODO(author): replace with your name, affiliation and ORCID before submission.
\author{R.\ [Surname]\thanks{Correspondence: \texttt{[email]}. Source code:
\url{https://github.com/dutch-rob/MyFeelsLike} (GPL-3.0-or-later).}}

\date{\today}

\begin{document}
\maketitle

\begin{abstract}
Operational ``feels like'' indices --- apparent temperature, heat index, wind
chill, humidex, UTCI --- are population constructs. They are calibrated to a
reference person with a reference clothing level, metabolic rate and
physiology, and they report a single number for everyone standing in the same
weather. Individuals depart from that reference substantially and
systematically: the same 12\,\textdegree C morning is pleasant to one person
and cold to another, and the disagreement is stable over time rather than
random. We describe MyFeelsLike, a smartphone application that estimates a
\emph{personal} apparent-temperature scale by regressing a user's own comfort
ratings on the weather that was present when each rating was made, and then
applies the fitted model to hourly forecasts. The contributions are (i) a
rating instrument that elicits a continuous comfort score without exposing
temperature units, thereby avoiding anchoring on the very quantity being
modelled; (ii) an estimation procedure designed for the small-$n$ regime that
a single user can realistically supply, combining a hard feature budget,
corrected-AIC forward selection and a set of physical plausibility screens that
reject coefficient sets implying an implausible response to temperature, wind
or activity; (iii) a multi-part reliability measure that widens or narrows the
displayed band according to statistical leverage and to how far a forecast
falls outside the range of conditions the user has actually rated; and (iv) an
evaluation protocol for comparing personalized against population indices. We
specify the method in full and state the evaluation design; empirical results
from the deployed user base are reported separately.
\end{abstract}

\section{Introduction}

Every weather service reports a ``feels like'' temperature, and every such
number is an average over people. \citet{steadman1984}, building on
\citet{steadman1979}, defines apparent temperature for an adult of specified height, weight and clothing walking at a
specified speed; the heat index, humidex \citep{masterton1979} and wind chill
\citep{osczevski2005} make comparable commitments. The Universal Thermal
Climate Index \citep{jendritzky2012} is derived from a multi-node
thermoregulation model of a reference person. These indices succeed at what
they were built for --- issuing warnings calibrated to a population --- and
they are not intended to describe any particular person.

The residual variation between people is, however, large enough to matter for
ordinary decisions: what to wear, whether to cycle, when to run. Indoor thermal
comfort research has long recognized this. \citet{fanger1970} predicts a mean
vote (PMV) and an expected percentage dissatisfied (PPD), the latter
acknowledging that a substantial fraction of occupants disagree with the mean
even at its optimum. Adaptive comfort models \citep{dedear1998} explain part of
the gap by context and acclimatization. More recently, \emph{personal comfort
models} \citep{kim2018} abandon the population target and fit a separate model
per occupant from that occupant's own feedback, typically in an instrumented
building with an HVAC system to control.

MyFeelsLike asks whether the personal-comfort-model idea transfers outdoors,
where the covariates cannot be manipulated, the sampling is opportunistic, the
label is supplied by an untrained rater on a phone, and the useful output is a
forecast rather than a setpoint. The setting imposes three constraints that
shape the entire design:

\begin{enumerate}
\item \textbf{Very small $n$.} An engaged user supplies tens of ratings, not
      thousands. The model must be useful at $n = 10$ and must not overfit at
      $n = 40$.
\item \textbf{Confounded, self-selected covariates.} People rate the weather
      they are in. Temperature, humidity and sunshine covary seasonally, and
      the user chooses when to rate, so the design matrix is neither balanced
      nor randomized.
\item \textbf{A label that must not be anchored.} If the instrument shows a
      temperature, the rating collapses toward that temperature. The scale must
      therefore carry no units.
\end{enumerate}

Section~\ref{sec:ratings} describes the rating instrument,
Section~\ref{sec:features} the feature construction, Section~\ref{sec:fitting}
estimation and model selection, Section~\ref{sec:forecast} application to
forecasts, Section~\ref{sec:reliability} the reliability measure,
Section~\ref{sec:presentation} the visual encoding and its accessibility
variants, and Section~\ref{sec:evaluation} the evaluation protocol.

\section{Collecting ratings}
\label{sec:ratings}

\subsection{The instrument}

A rating is elicited on a continuous, unitless scale from 0 to 1000, presented
as a vertical colour column that the user scrolls until the colour under a
fixed indicator matches how the present weather feels. The scale runs white
(coldest) through blue, green, yellow and red to dark red (hottest), with
anchor positions at 0, 250, 500, 690, 880 and 1000 respectively; intermediate
scores interpolate linearly in RGB between neighbouring anchors. Green sits at
the midpoint by design, so that ``comfortable'' is the centre of the scale
rather than a point along a cold-to-hot ramp.

Two properties of this instrument are deliberate. First, \emph{no temperature
is displayed}: the score is the response variable, and showing a temperature
alongside it would invite the user to report the thermometer rather than their
own sensation. Second, the scale is continuous rather than the discrete
seven-point ASHRAE thermal sensation scale. A continuous score avoids the
floor and ceiling effects that a short ordinal scale produces at climatic
extremes, and it supplies more information per observation --- which matters
when observations are scarce. For users who prefer words, or who cannot use
the colour scale, the same 0--1000 range is presented as six named categories
(\emph{very cold, cold, mild, warm, very warm, hot}) placed at the colour
anchors, with linear transitions between them; the recorded quantity is
identical.

\subsection{Context supplied by the user}

Three contextual variables are recorded with each rating, because they change
the sensation without changing the weather:

\begin{itemize}
\item \textbf{Activity} (0 not active, 1 light, 2 moderate, 3 vigorous), a
      coarse proxy for metabolic rate.
\item \textbf{Dress} ($-2$ dressed for very cold $\ldots$ $0$ dressed
      appropriately $\ldots$ $+2$ dressed for very warm), a coarse proxy for
      clothing insulation relative to conditions.
\item \textbf{Sun exposure} ($+1$ full sun, $0$ partial, $-1$ shade).
\end{itemize}

These are deliberately coarse. Asking for clo values or MET rates would raise
the cost of a rating and lower the response rate, and the resulting precision
would be spurious given the rest of the measurement chain.

\subsection{Covariates captured automatically}

At the moment of rating, the application records the concurrent weather for the
user's location from the platform weather service: dry-bulb temperature,
service-reported apparent temperature, dew point, relative humidity, station
pressure, wind speed, precipitation probability and amount, total cloud cover
and its low, medium and high components, UV index, and a daylight flag. Wet-bulb
temperature is derived from temperature and relative humidity using
\citet{stull2011}. Storing the covariates at rating time, rather than
re-fetching them later, is what makes the observation immutable: a subsequent
revision of the forecast archive cannot silently change a training row.

\subsection{Eligibility for fitting}

No model is fitted until two conditions hold: at least five ratings exist, and
the spread of the recorded scores (maximum minus minimum) is at least 50 points
--- five per cent of the scale. The second condition matters more than the
first. A user who has rated ten mild mornings at 500 supplies almost no
information about their response to temperature, and a regression fitted to
that data would extrapolate wildly. Until both hold, the application shows the
population apparent temperature and explains which condition is unmet.

\section{Feature construction}
\label{sec:features}

The candidate feature set contains 28 terms in four groups.

\paragraph{Direct meteorological terms.} Apparent temperature, dry-bulb
temperature, wet-bulb temperature, dew point, relative humidity, station
pressure, wind speed, precipitation probability and amount, total, low, medium
and high cloud cover, UV index, and the daylight indicator.

\paragraph{Derived differences.} Apparent minus dry-bulb (the service's own
wind and humidity correction), dry-bulb minus wet-bulb (wet-bulb depression),
and wet-bulb minus dew point. These are included because they isolate physical
effects that the raw variables confound: two days at the same temperature with
different wet-bulb depressions are physiologically different days.

\paragraph{Contextual terms.} Activity, dress and sun exposure as recorded
above.

\paragraph{Piecewise-linear hinges and interactions.} Human thermal response is
not linear in temperature; the same one-degree change matters more near the
extremes than in the middle. Rather than fitting splines, which the sample size
will not support, the candidate set includes five hinge terms of the form
$\max(0, x - c)$ or $\max(0, c - x)$ with fixed knots chosen on physical
grounds: cold amplification below 10\,\textdegree C, warm onset above
18\,\textdegree C, heat amplification above 26\,\textdegree C, a wind threshold
at 15\,kph, and a UV threshold at 4. Three interaction terms (apparent
temperature with humidity, with UV, and with activity) capture the familiar
effects that humidity matters more when it is hot and that exertion matters
more when it is warm.

Hinges become candidates only at $n \geq 25$ and interactions only at
$n \geq 40$. A hinge estimated from five observations describes noise.

All features are standardized before fitting, so estimated coefficients are
expressed per standard deviation rather than per physical unit. This is
essential to keep in mind when interpreting them, and the plausibility screens
of Section~\ref{sec:plausibility} therefore work in physical units by probing
the fitted model rather than by inspecting coefficients directly.

\section{Estimating the model}
\label{sec:fitting}

\subsection{Objective and estimator}

Let $y_i \in [0, 1000]$ be the $i$-th comfort score and $\mathbf{x}_i$ the
standardized feature vector at that observation. The model is ordinary least
squares with an intercept,
\begin{equation}
y_i = \beta_0 + \sum_{j \in S} \beta_j x_{ij} + \varepsilon_i ,
\end{equation}
where $S$ is the selected feature subset. Coefficients are obtained by Cholesky
solution of the normal equations, and $(\mathbf{X}^\top\mathbf{X})^{-1}$ is
retained --- it is needed later for leverage.

\subsection{Selection}

Selection is forward stepwise on the corrected Akaike information criterion
\citep{hurvich1989, burnham2002},
\begin{equation}
\mathrm{AIC_c} = n \log\!\left(\frac{\mathrm{RSS}}{n}\right) + 2p +
\frac{2p(p+1)}{n - p - 1},
\end{equation}
with $p$ the number of estimated parameters. The correction term dominates in
exactly the regime of interest, and without it the criterion would happily add
features to a ten-observation fit.

Three rules constrain the search:

\begin{enumerate}
\item \textbf{Feature budget.} At most $\lfloor n/5 \rfloor$ features may enter,
      independently of what $\mathrm{AIC_c}$ prefers. Five observations per
      estimated coefficient is a conventional floor for regression, and here it
      is enforced as a hard cap rather than a guideline.
\item \textbf{Margin.} A feature enters only if it improves $\mathrm{AIC_c}$ by
      at least 2 over the best model without it; otherwise the search stops.
      The search begins from the intercept-only model, so the margin applies to
      the first feature as well as to later ones.
\item \textbf{Per-step eligibility.} Plausibility is checked at each step, on
      the model that would result, not once at the end.
\end{enumerate}

\subsection{Plausibility screens}
\label{sec:plausibility}

A model can fit the training ratings well and still be absurd --- the failure
that motivated this section was a real user seeing a forecast at the coldest
end of the scale during a heat wave, produced by a well-fitting model whose
temperature response had the wrong sign. Because coefficients are standardized
and features are correlated, the sign of any individual coefficient is not
diagnostic. The screens therefore \emph{probe} the candidate model: they
evaluate its predictions at the user's median conditions and at perturbations
of them, and reject the candidate if the implied physical response is not
credible.

A candidate feature is ineligible at a given step if adding it would imply any
of the following, where the response is measured at the user's median
conditions:

\begin{itemize}
\item a temperature response outside $[0, 60]$ score points per
      \textdegree C, evaluated by moving all temperature-like variables
      together by $\pm 3$\,\textdegree C (a negative response means the model
      claims the user feels colder as it gets warmer; 60 points per degree
      means the full scale is traversed in under 17 degrees);
\item a warmer reading from an additional 10\,kph of wind, beyond a tolerance
      of 25 score points;
\item a warmer reading from one additional level of activity, beyond the same
      tolerance;
\item any non-finite coefficient, mean or standard deviation, or a
      standard deviation that is not strictly positive.
\end{itemize}

A feature is also ineligible unless the observed spread of its underlying
physical variable reaches a per-variable minimum (for example 3\,\textdegree C
for temperature-like variables, 5\,kph for wind, 0.1 for relative humidity).
Fitting a wind coefficient from ratings that all occurred in still air is not
informative regardless of what the criterion says.

The wind and activity screens deserve comment, since they encode a physical
prior rather than a statistical one. It is possible for a person to genuinely
feel warmer in wind, and the tolerance of 25 points permits a modest effect in
that direction; what the screen rejects is a large one, which in this data is
overwhelmingly more likely to be confounding --- windy days in the user's
sample happening also to be warm days --- than physiology.

\section{Applying the model to forecasts}
\label{sec:forecast}

The fitted model is applied to each hour of the hourly forecast, and to each
day of the ten-day forecast, producing a personal score for every point.
Because the contextual variables are not forecast quantities, the user selects
a \emph{scenario} --- an activity level, a dress level and a sun exposure ---
which is held fixed across the forecast horizon. The default scenario is the
user's modal rating context.

When the fitted model contains a sun-exposure term, two predictions are
produced per hour: one in shade and one in sun. After dark the two are
deliberately collapsed onto the shade value, since ``in sun'' is not a state
that exists at night; the pair is displayed either as two bands or as a
gradient within one band.

Predictions are clamped to $[0, 1000]$ for display. Clamping is a display
decision, not a modelling one: the underlying prediction is retained for the
reliability calculation, where an out-of-range prediction is itself evidence
that the model is being extrapolated.

\section{Reliability}
\label{sec:reliability}

A personalized model built from tens of observations will be asked to predict
conditions the user has never rated. Communicating that honestly is as
important as the point prediction, and a single confidence interval is a poor
vehicle: the reasons a prediction is untrustworthy are heterogeneous, and the
user's response to each is different.

The application computes four widths, each in $[0.15, 1]$, and takes their
minimum. The floor of 0.15 means the displayed band never vanishes entirely; a
sliver still communicates a value while making the loss of confidence obvious.

\paragraph{Leverage.} With $h = \mathbf{x}^\top(\mathbf{X}^\top\mathbf{X})^{-1}
\mathbf{x}$ the leverage of the forecast point in the standardized design,
$m$ the number of parameters and $n$ the number of ratings, the width is $1$
for $h \leq 2m/n$, $0$ for $h \geq 3m/n$, and linear between. This is the
classical high-leverage rule of thumb used as a continuous fade rather than a
flag.

\paragraph{In-model range.} For each selected feature, the range of values
observed in training is recorded. A forecast value inside that range gives
width 1; one more than a full range-length outside it gives 0.15; linear
between. Leverage alone does not capture this cleanly, because a point can have
modest leverage while still lying beyond every value the user has rated.

\paragraph{Not-in-model range.} Variables that are \emph{not} in the model are
also checked, with a per-variable tolerance factor. A model that omits wind
because the user only ever rated calm days should not be silently trusted on a
forecast gale. Tolerances vary by variable: temperature-like variables use
1.5 range-lengths, humidity, wind and UV use 2, and precipitation and cloud
variables use 4, reflecting how strongly each is expected to matter when it
falls outside the sampled range.

\paragraph{Predicted-score range.} A prediction far outside the range of scores
the user has actually given, or a non-finite prediction, narrows the band to
the floor.

The resulting width is encoded geometrically rather than by transparency
wherever possible: the hourly band's thickness about its centre line, and the
ten-day heatmap's cell width within its day column. Where geometry is not
available --- the sun/shade gradient, which needs full height to be legible ---
opacity is used with a floor of 0.25.

\section{Presentation and accessibility}
\label{sec:presentation}

The comfort scale is colour by default, which makes it inaccessible to a
non-trivial minority of users. Roughly 8\% of men and 0.5\% of women of
Northern European descent have some form of colour vision deficiency, the
large majority red--green (deuteranomaly and protanomaly). The application
therefore offers three encodings of the same underlying score:

\begin{enumerate}
\item \textbf{Full colour}, as described in Section~\ref{sec:ratings}.
\item \textbf{Red-free}, which keeps every anchor up to pure yellow in place
      and replaces the red end with yellow darkening to black. What separates
      the hot end from the middle is then lightness rather than hue, so
      red--green confusion cannot make a hot hour read as a mild one. Because
      the shared anchors do not move, a rating made under one encoding means
      the same under the other.
\item \textbf{No colour}, in which hue carries nothing. Comfort is shown either
      as a line graph against the six named categories, or as the same bands
      rendered in grayscale; the weather curves are distinguished by line
      pattern (solid, dashed, dotted, dash-dot) rather than colour; and the
      rating scale states both the category name and the numeric score.
\end{enumerate}

This matters methodologically and not only ethically: the rating instrument
\emph{is} the measurement device, and a user who cannot read the scale
reliably supplies noisier labels.

\section{Evaluation protocol}
\label{sec:evaluation}

We state the evaluation design here; measured results are not included in this
draft and will be reported once the deployed data set is of sufficient size.

\subsection{Question}

Does a personalized model predict an individual's comfort ratings better than
population indices do? The comparison is against three baselines: (i) dry-bulb
temperature, (ii) the weather service's apparent temperature, mapped
monotonically onto the score scale, and (iii) a pooled model fitted across all
users, which isolates the value of \emph{personalization} from the value of
\emph{regression}.

\subsection{Design}

The primary estimate is leave-one-out cross-validation within user, which
matches the small-$n$ regime and uses every observation. Because ratings are
temporally autocorrelated --- consecutive ratings are often the same day ---
LOOCV will be optimistic, and a blocked variant leaving out whole days is
reported alongside it. The strongest evidence is prospective: predictions are
recorded before the rating is made, and scored when it arrives.

\subsection{Metrics}

Mean absolute error in score units is the headline metric, being interpretable
on the instrument the user actually sees. Skill relative to each baseline is
reported as the proportional reduction in MAE. Two further quantities matter
specifically here:

\begin{itemize}
\item \textbf{Learning curve.} MAE as a function of $n$, which answers the
      practical question of how many ratings a user must supply before the
      model is worth showing.
\item \textbf{Reliability calibration.} Whether the width described in
      Section~\ref{sec:reliability} is honest --- that is, whether absolute
      error is in fact larger where the band is narrower. A reliability signal
      that does not track error is worse than none, since it invites misplaced
      confidence.
\end{itemize}

\subsection{Between-person variation}

The application supports comparing forecasts between consenting users, which
provides a direct estimate of the quantity that motivates the work: the spread
of personal scores across people under identical weather. Reporting that
spread, and its stability over time, tests the premise that individual
differences are systematic rather than noise.

\section{Limitations}

\paragraph{Opportunistic sampling.} Users rate when they choose to, producing a
design matrix that is neither balanced nor randomized. Seasonal confounding is
the most serious consequence: over a single season, temperature, humidity, UV
and daylight move together, and no amount of model selection can separate
effects that never vary independently. The plausibility screens mitigate the
worst outcomes but do not solve the underlying identification problem.

\paragraph{Coarse context.} Activity, dress and sun are three-to-five level
ordinal proxies for metabolic rate, clothing insulation and radiant load.

\paragraph{Single-rater labels.} There is no way to estimate the measurement
error of the instrument from a single rater, since a rating cannot be repeated
under identical conditions. Test--retest reliability would require deliberate
duplicate ratings within a short window.

\paragraph{Adaptation.} Thermal preference drifts with acclimatization
\citep{dedear1998}. The model as described weights all ratings equally
regardless of age; a user's ratings from the previous winter may not describe
this one. Recency weighting is an obvious extension and is not yet implemented.

\paragraph{Forecast error.} The personal model is applied to forecast
covariates, so its errors compound with the underlying forecast's. The
reliability measure captures uncertainty in the personal model only.

\section{Availability and privacy}

Ratings, the fitted model and the derived forecasts remain on the user's
device and in their private cloud account; nothing is transmitted to the
authors by default. Sharing a model with another user is explicit and
per-recipient, and is implemented as a self-contained snapshot rather than a
live feed. Source code is available at
\url{https://github.com/dutch-rob/MyFeelsLike} under GPL-3.0-or-later.

\section{Conclusion}

Population apparent-temperature indices answer a question no individual asked:
how does the average person feel? MyFeelsLike substitutes the question each
user actually has, and shows that the substitution is tractable on a phone with
a few tens of self-supplied ratings --- provided the estimator is built for
that regime rather than borrowed from one with abundant data, and provided the
interface is honest about the many forecasts for which the answer is not yet
known.

\begin{thebibliography}{99}

\bibitem[Burnham and Anderson(2002)]{burnham2002}
Burnham, K.~P. and Anderson, D.~R. (2002).
\newblock \emph{Model Selection and Multimodel Inference: A Practical
  Information-Theoretic Approach}, 2nd edition.
\newblock Springer, New York.

\bibitem[de Dear and Brager(1998)]{dedear1998}
de Dear, R.~J. and Brager, G.~S. (1998).
\newblock Developing an adaptive model of thermal comfort and preference.
\newblock \emph{ASHRAE Transactions}, 104(1):145--167.

\bibitem[Fanger(1970)]{fanger1970}
Fanger, P.~O. (1970).
\newblock \emph{Thermal Comfort: Analysis and Applications in Environmental
  Engineering}.
\newblock Danish Technical Press, Copenhagen.

\bibitem[Hurvich and Tsai(1989)]{hurvich1989}
Hurvich, C.~M. and Tsai, C.-L. (1989).
\newblock Regression and time series model selection in small samples.
\newblock \emph{Biometrika}, 76(2):297--307.

\bibitem[Jendritzky et~al.(2012)]{jendritzky2012}
Jendritzky, G., de Dear, R., and Havenith, G. (2012).
\newblock UTCI---why another thermal index?
\newblock \emph{International Journal of Biometeorology}, 56(3):421--428.

\bibitem[Kim et~al.(2018)]{kim2018}
Kim, J., Zhou, Y., Schiavon, S., Raftery, P., and Brager, G. (2018).
\newblock Personal comfort models: Predicting individual thermal preference
  using occupant heating and cooling behavior.
\newblock \emph{Building and Environment}, 129:96--106.

\bibitem[Masterton and Richardson(1979)]{masterton1979}
Masterton, J.~M. and Richardson, F.~A. (1979).
\newblock \emph{Humidex: A Method of Quantifying Human Discomfort Due to
  Excessive Heat and Humidity}.
\newblock Environment Canada, Atmospheric Environment Service, Downsview,
  Ontario.

\bibitem[Osczevski and Bluestein(2005)]{osczevski2005}
Osczevski, R. and Bluestein, M. (2005).
\newblock The new wind chill equivalent temperature chart.
\newblock \emph{Bulletin of the American Meteorological Society},
  86(10):1453--1458.

\bibitem[Steadman(1979)]{steadman1979}
Steadman, R.~G. (1979).
\newblock The assessment of sultriness. Part I: A temperature-humidity index
  based on human physiology and clothing science.
\newblock \emph{Journal of Applied Meteorology}, 18(7):861--873.

\bibitem[Steadman(1984)]{steadman1984}
Steadman, R.~G. (1984).
\newblock A universal scale of apparent temperature.
\newblock \emph{Journal of Climate and Applied Meteorology}, 23(12):1674--1687.

\bibitem[Stull(2011)]{stull2011}
Stull, R. (2011).
\newblock Wet-bulb temperature from relative humidity and air temperature.
\newblock \emph{Journal of Applied Meteorology and Climatology},
  50(11):2267--2269.

\end{thebibliography}


\end{document}
